\note{In UC3, the authors test a clear design objective for multipath transport: using two paths should never perform worse than using just one of those paths alone.\par \par They first evaluated \textbf{DChannel}, which combines a high-bandwidth path (HBP) and a low-latency path (LLP) in parallel while steering packets transparently below the transport layer. Here, the reference uses only the high-bandwidth path, the target uses DChannel over both paths, and the metric is flow completion time (FCT) for downloading a 20 MB file. Although prior work saw no issues when running loss-based CUBIC over DChannel, AdvNet found an environment—where the HBP has a 120 ms one-way delay and the LLP has only a 3 ms one-way delay—in which DChannel's throughput drops to one-third of using the HBP alone.\par \par The root cause lies in DChannel's packet-steering heuristic, which sends small control packets (such as TCP handshakes) over the fast LLP and large data packets over the slower HBP. Because the initial handshake completes over the 3 ms path, TCP estimates a tiny RTT and sets its retransmission timeout (RTO) to the Linux minimum of roughly 200 ms. When large data packets subsequently take the 120 ms one-way HBP, their round-trip and queuing delay exceeds the 200 ms RTO before ACKs can return—triggering spurious timeouts and rate reductions.
\par\smallskip Source: AdvNet paper §4.4.3}
